\scanpage{49}
\uline{Ainsi, on trouve que}, \inserted{\uline{posant} $(\unsure{a|*})$,} \uline{les complexes} $g^*(K^{\cdot})$, \uline{pour les différentes sections} $g$ \uline{de} $M$ \uline{sur} $S$ \uline{passant par} $t_0$, \uline{sont can. isom entre eux} \uline{(au sens des cat. dérivées)} \uline{donc les} $g^*(\mathscr H^i)$ (\uline{où} $\mathscr H^i=H^i(K^{\cdot})\inserted{=\mathbb R^{\unsure{i|j|\nu}}f_*(\Omega^{\cdot}_{Y/M})}$) \uline{sont \unsure{canoniqt|canoniqmt|canoniquement} isom entre eux}, \deleted{\unsure{pour}} \uline{mod torsion}. La donnée d'un tel système transitif d'isomorphismes peut être considéré comme une \uline{structure \unsure{bourbaquiste|bourbakiste} sur} $K^{\unsure{*|\cdot}}=\mathbb R f_*(\deleted{\unsure{Y}}\underline{\Omega}^{\cdot}_{Y/M})$ [en \uuline{le point} $t_0$, relativement à la base $S$], ou sur les $\mathscr H^i=\mathbb R^i f_*(\Omega^{\cdot}_{Y/M})$, \deleted{Une telle structure est \unsure{présente|présentée} p.ex. autrement} \inserted{d'espèce voisine à une \uline{structure de}} \uline{descente}. Une structure de \deleted{\unsure{cet|ce}} ce dernier type, \inserted{pour un Module $\underline E$ disons,} : supposons qu'elles sont effectives, i.e. définissent un isom
\[
\underline E\deleted{\mathscr H^{\unsure{i}}}\simeq h^{\unsure{+|*}}(\underline E_0)
\]
($\underline E_0$ un Module sur $S$], définit \unsure{un|une}
% archive sheet number: 48

